% ============================================================
%  DÉROULEMENT DU PROJET — deroulement.tex
% ============================================================
\section{Déroulement du projet}
\label{sec:deroulement}

Cette section retrace l'organisation du travail d'équipe et les étapes successives
ayant conduit à l'application décrite dans la section~\ref{sec:realisation}. Elle
présente la méthode de collaboration adoptée, la répartition des responsabilités,
les outils de développement utilisés, et les principaux obstacles surmontés au
cours du projet.

% ------------------------------------------------------------
\subsection{Organisation et répartition du travail}
\label{subsec:organisation_travail}

L'équipe est composée de trois développeurs~: \textsc{Feraoun} Mohamed~Amine,
\textsc{Anurajan} Thenuxshan et \textsc{Ouardani} Nidhal. Afin de tirer parti de
la modularité de l'architecture décrite à la section~\ref{subsec:archi_globale},
la répartition des tâches a suivi un découpage fonctionnel cohérent avec les
quatre modules du projet.

\paragraph{Répartition du travail.}
Feraoun Mohamed Amine a conçu l'architecture modulaire et le cœur du moteur de
simulation (classes Carte, ManageurBasique, règles de transformation). Il a
réalisé les tests unitaires et rédigé la première partie du rapport.

Ouardani Nidhal a mis en place le système de journalisation Log4j et a développé
l'algorithme de génération cohérente de la carte et la logique de déplacement
des événements mobiles. Il a contribué à la rédaction de la deuxième partie du
rapport.

Anurajan Thenuxshan a implémenté le module de configuration et la fenêtre
d'édition (PanelEdition), assurant l'interface de paramétrage des événements. Il a
apporté un support technique sur les composantes graphiques et participé aux tests
d'intégration.

Cette répartition n'a pas été strictement cloisonnée~: des révisions croisées du
code ont été systématiquement organisées pour garantir la cohérence globale du
projet et mutualiser les connaissances techniques entre les membres de l'équipe.

\paragraph{Méthode de travail collaborative.}
Un dépôt Git commun~\cite{chacon2014pro} (lien~: \url{https://github.com/MLX-JAY/Environnement}) a servi de base au travail collaboratif.
Chaque fonctionnalité a été développée sur une branche dédiée avant d'être
fusionnée vers la branche principale après revue de code. Des réunions d'équipe
hebdomadaires ont permis de synchroniser les avancées, d'identifier les dépendances
entre modules et d'ajuster le planning.

% ------------------------------------------------------------
\subsection{Phases de développement}
\label{subsec:phases}

Le développement s'est articulé en quatre phases successives dont les jalons sont
présentés dans le tableau~\ref{tab:planning}.

\begin{table}[htbp]
  \centering
  \caption{Phases de développement et jalons du projet}
  \label{tab:planning}
  \begin{tabular}{cl>{\raggedright\arraybackslash}p{8cm}}
    \toprule
    \textbf{Phase} & \textbf{Période} & \textbf{Objectifs et livrables} \\
    \midrule
    1 & Semaines 1--2  & Analyse des besoins, conception de l'architecture,
                         mise en place de l'environnement de développement
                         et du dépôt Git \\[0.4em]
    2 & Semaines 3--5  & Implémentation des classes de données (biomes,
                         événements, carte), module \texttt{config},
                         premières règles de transformation \\[0.4em]
    3 & Semaines 6--8  & Développement du moteur de simulation
                         (\texttt{ManageurBasique}), intégration des factories,
                         réalisation de l'interface graphique \\[0.4em]
    4 & Semaines 9--10 & Finalisation des différents concepts du projet
                         (comme le mode édition), corrections des anomalies
                         et bilan de fin ; rédaction du rapport \\
    \bottomrule
  \end{tabular}
\end{table}

% ------------------------------------------------------------
\subsection{Difficultés rencontrées et solutions adoptées}
\label{subsec:difficultes}

Plusieurs difficultés techniques ont été rencontrées au cours du développement.
Elles sont décrites dans la section~\ref{subsec:difficultes} sous la forme d'un couple problème/solution afin
d'en faciliter la compréhension.

\paragraph{Les statistiques.}
\textbf{Problème.}
Une première version des statistiques reposait sur des bulles visuelles se
remplissant proportionnellement à l'état du monde. Cette représentation s'est
révélée imprécise et insuffisante pour suivre l'évolution quantitative de
l'écosystème au fil des rounds.

\textbf{Solution.}
La bibliothèque JFreeChart a été intégrée pour remplacer ce dispositif. Elle
permet d'afficher des graphiques en courbes clairs alimentés en temps réel par les
données de simulation, offrant une lecture précise et rigoureuse des propriétés
des biomes au fil du temps.

\paragraph{Les transformations écologiques.}
\textbf{Problème.}
La gestion dynamique des transformations de biomes s'est avérée complexe à
implémenter~: un premier prototype fondé sur des conditions imbriquées dans le
\texttt{ManageurBasique} était difficile à maintenir et ne tenait pas compte des
spécificités de chaque biome.

\textbf{Solution.}
Un système de règles de transformation a été créé, matérialisé par l'interface
\texttt{RegleTransformation}. Chaque règle est liée aux statistiques propres du
biome concerné et encapsule sa logique de manière autonome. Le moteur se contente
d'itérer sur la liste des règles, ce qui rend l'ensemble extensible et maintenable.

\paragraph{Dynamique des biomes.}
\textbf{Problème.}
Dans une version initiale, la seule différence entre les biomes était d'ordre
visuel (couleur et texture). Tous les biomes recevaient les mêmes types
d'événements, ce qui rendait la simulation peu réaliste et les identités de biomes
interchangeables.

\textbf{Solution.}
Le patron Visiteur (\texttt{GenerateurEvenementVisitor}) a été mis en place pour
que les événements générés dépendent du type de biome visité. Un biome Mer génère
ainsi des événements de type Pluie, tandis qu'un biome Désert favorise des
événements de Zéphyr ou de Tornade. Cette logique renforce l'identité unique de
chaque biome.

\paragraph{Création des événements.}
\textbf{Problème.}
L'apparition d'un événement sur une seule case était peu lisible et manquait de
cohérence visuelle. Un événement de type Orage occupant une cellule unique avait le
même impact visuel qu'une simple Pluie, rendant la lecture de la grille confuse.

\textbf{Solution.}
Un système de formations spatiales a été implémenté directement dans
l'\texttt{EvenementFactory}~: chaque événement est désormais créé selon une
formation géométrique définie (ligne, carré, croix, etc.), via des méthodes de
création dédiées. Ce mécanisme apporte un boost visuel et logique significatif~:
chaque type d'événement possède une empreinte reconnaissable sur la grille, ce qui
renforce la cohérence globale de l'interface.
